\section{经典线性正算子与 Korovkin 定理}

Weierstrass 第一逼近定理和第二逼近定理分别肯定了闭区间上的连续函数可用多项式一致逼近，周期连续函数可用三角多项式一致逼近。然而，这两个定理本身并未提供统一的构造方法。本章将从具体的构造性算子入手，给出 Bernstein、Baskakov、Mirakyan 以及 Vallée-Poussin 算子的直接证明，随后引入 Korovkin 定理，将这些算子的收敛性纳入一个简洁的统一框架之中。

\subsection{Weierstrass 第一逼近定理与 Bernstein 多项式}

Weierstrass 第一逼近定理断言：闭区间上的任何连续函数都能用一列多项式一致逼近。用数学语言描述，就是
\[
\forall \epsilon > 0,\ \exists N,\ \forall n \geq N,\ \forall x \in [a, b],\ |f(x) - p_n(x)| < \epsilon,
\]
其中 $\{p_n(x)\}$ 是一列多项式。下面给出具体证明。

我们先考虑 $f(x)$ 在闭区间 $[0,1]$ 上的情形。使用 Bernstein 多项式来构造逼近函数：
\[
B_n^f(x) = \sum_{k=0}^{n} f\left(\frac{k}{n}\right) \binom{n}{k} x^k (1-x)^{n-k}.
\]
构造的动机与概率论中的二项分布密切相关：进行 $n$ 次独立试验，每次成功概率为 $p$，成功 $k$ 次的概率为
\[
P(X = k) = \binom{n}{k} p^k (1-p)^{n-k}.
\]
当 $k$ 接近 $np$ 时概率最大。因此 $B_n^f(x)$ 相当于对 $f$ 在 $k/n\approx x$ 处的取值进行加权平均，使得 $B_n^f(x)$ 能更好地逼近 $f(x)$。

为证明一致收敛，需要以下组合恒等式。

\begin{lemma}
\[
\sum_{k=0}^{n} \binom{n}{k} u^k v^{n-k} = (u + v)^n.
\]
\end{lemma}
\begin{lemma}
\[
\sum_{k=0}^{n} k \binom{n}{k} u^k v^{n-k} = nu(u + v)^{n-1}.
\]
\end{lemma}
\begin{lemma}
\[
\sum_{k=0}^{n} k^2 \binom{n}{k} u^k v^{n-k} = nu(u+v)^{n-1} + n(n-1)u^2(u+v)^{n-2}.
\]
\end{lemma}

这些可由二项式定理对 $u$ 求导得到。特别地，取 $u=v=1$ 可得推论：
\[
\sum_{k=0}^{n} k \binom{n}{k} = n \cdot 2^{n-1},\qquad
\sum_{k=0}^{n} k^2 \binom{n}{k} = n(n+1)2^{n-2}.
\]

代入 $u=x,\ v=1-x$，由引理1得
\[
\sum_{k=0}^{n} \binom{n}{k} x^k (1-x)^{n-k} = 1,
\]
并可计算出关键的方差型求和：
\begin{align*}
\sum_{k=0}^{n} (k-nx)^2 \binom{n}{k} x^k (1-x)^{n-k}
&= \sum_{k=0}^{n} k^2 \binom{n}{k} x^k (1-x)^{n-k} \\
&\quad - 2nx\sum_{k=0}^{n} k\binom{n}{k}x^k(1-x)^{n-k} + n^2x^2\sum_{k=0}^{n}\binom{n}{k}x^k(1-x)^{n-k} \\
&= [nx + n(n-1)x^2] - 2nx\cdot(nx) + n^2x^2 \\
&= nx(1-x).
\end{align*}

现在证明 $B_n^f(x)$ 一致收敛于 $f(x)$。

\begin{proof}
由于 $f$ 在 $[0,1]$ 上连续，故一致连续：
\[
\forall \epsilon>0,\ \exists\delta>0,\ \forall x,y\in[0,1],\ |x-y|<\delta \Rightarrow |f(x)-f(y)|<\epsilon.
\]
利用 $\sum_{k=0}^{n} \binom{n}{k} x^k (1-x)^{n-k} = 1$，可将 $f(x)$ 改写为
\[
f(x) = \sum_{k=0}^{n} f(x) \binom{n}{k} x^k (1-x)^{n-k},
\]
于是
\[
B_n^f(x)-f(x) = \sum_{k=0}^{n} \left(f\left(\frac{k}{n}\right)-f(x)\right) \binom{n}{k} x^k (1-x)^{n-k}.
\]
按 $|k/n-x|<\delta$ 和 $\ge\delta$ 拆分估计。

当 $|k/n-x|<\delta$ 时，$|f(k/n)-f(x)|<\epsilon$，故
\[
\left|\sum_{|k/n-x|<\delta} \cdots \right|
\le \sum_{|k/n-x|<\delta} \epsilon \binom{n}{k} x^k (1-x)^{n-k}
\le \epsilon \sum_{k=0}^{n} \binom{n}{k} x^k (1-x)^{n-k} = \epsilon.
\]

当 $|k/n-x|\ge\delta$ 时，设 $|f(x)|\le M$，则 $|f(k/n)-f(x)|\le 2M$。利用 $(k-nx)^2 \ge n^2\delta^2$ 放大：
\[
\left|\sum_{|k/n-x|\ge\delta} \cdots \right|
\le \frac{2M}{n^2\delta^2} \sum_{k=0}^{n} (k-nx)^2 \binom{n}{k} x^k (1-x)^{n-k}
= \frac{2M}{n^2\delta^2}\cdot nx(1-x).
\]
因为 $x(1-x)\le 1/4$，上式 $\le \frac{M}{2n\delta^2}$。

综合两部分得
\[
|B_n^f(x)-f(x)| < \epsilon + \frac{M}{2n\delta^2}.
\]
取 $N > \frac{M}{2\epsilon\delta^2}$，则当 $n>N$ 时 $|B_n^f(x)-f(x)| < 2\epsilon$，故 $B_n^f(x)$ 在 $[0,1]$ 上一致收敛于 $f(x)$。
\end{proof}

对于一般闭区间 $[a,b]$ 上的连续函数 $f(x)$，通过线性变换
\[
g(t) = f(a+(b-a)t),\quad t\in[0,1],
\]
将问题化归到 $[0,1]$ 上。$g(t)$ 仍连续，存在 Bernstein 多项式 $B_n^g(t)$ 一致收敛于 $g(t)$。将 $t=\frac{x-a}{b-a}$ 代回，得到
\[
B_n^f(x) = B_n^g\!\left( \frac{x-a}{b-a} \right)
= \sum_{k=0}^{n} f\!\left( a+\frac{k}{n}(b-a) \right)
\binom{n}{k} \left( \frac{x-a}{b-a} \right)^{\!k}
\left( 1-\frac{x-a}{b-a} \right)^{\!n-k},
\]
它在 $[a,b]$ 上一致收敛于 $f(x)$。这就完成了 Weierstrass 第一逼近定理的完整证明。


\subsection{Fejér 算子与 Weierstrass 第二逼近定理}

Weierstrass 第二逼近定理断言：任何以 $2\pi$ 为周期的连续函数都能用一列三角多项式一致逼近。
Fejér 给出了一个构造性的证明，其核心是对傅里叶级数的部分和取算术平均。

设 $f\in C_{2\pi}$，其傅里叶级数的前 $n+1$ 个部分和为 $S_0(f,x),S_1(f,x),\dots ,S_n(f,x)$。
定义
\[
\sigma_n(f,x)=\frac{S_0(f,x)+S_1(f,x)+\cdots +S_n(f,x)}{n+1},
\]
称为 \textbf{Fejér 算子}。每个 $\sigma_n(f,x)$ 都是次数不超过 $n$ 的三角多项式。

由傅里叶级数理论，部分和可表示为 Dirichlet 核的卷积：
\[
S_k(f,x)=\frac1\pi\int_{-\pi}^{\pi}f(x-t)D_k(t)\,dt,\qquad 
D_k(t)=\frac{\sin\bigl((k+\frac12)t\bigr)}{2\sin\frac{t}{2}}.
\]
代入平均得
\[
\sigma_n(f,x)=\frac1\pi\int_{-\pi}^{\pi}f(x-t)\mathcal{F}_n(t)\,dt,
\]
其中
\[
\mathcal{F}_n(t)=\frac1{n+1}\sum_{k=0}^{n}D_k(t)
=\frac{1}{2(n+1)}\frac{\sin^2\!\bigl(\frac{n+1}{2}t\bigr)}{\sin^2\frac{t}{2}}
\]
称为 \textbf{Fejér 核}。它有两个基本性质：
\[
\mathcal{F}_n(t)\ge 0,\qquad 
\frac1\pi\int_{-\pi}^{\pi}\mathcal{F}_n(t)\,dt=1.
\]
非负性来自表达式中分子分母均非负，归一性则因为每个 $D_k$ 的积分为 $\pi$。

下面证明 $\sigma_n(f,x)$ 一致收敛于 $f(x)$，证明结构与 Bernstein 多项式的情形完全平行。

\begin{proof}
由于 $f\in C_{2\pi}$，它在 $\mathbb{R}$ 上一致连续：
\[
\forall\varepsilon>0,\ \exists\delta>0,\ \forall x,y:\ |x-y|<\delta\Rightarrow |f(x)-f(y)|<\varepsilon.
\]
利用归一化条件，将 $f(x)$ 写成积分形式：
\[
f(x)=\frac1\pi\int_{-\pi}^{\pi}f(x)\mathcal{F}_n(t)\,dt,
\]
于是
\[
\sigma_n(f,x)-f(x)=\frac1\pi\int_{-\pi}^{\pi}\bigl[f(x-t)-f(x)\bigr]\mathcal{F}_n(t)\,dt.
\]
将积分区间按 $|t|<\delta$ 和 $\delta\le|t|\le\pi$ 拆成两部分。

\textbf{第一部分} $|t|<\delta$。此时 $|f(x-t)-f(x)|<\varepsilon$，结合 $\mathcal{F}_n(t)\ge0$，
\[
\frac1\pi\biggl|\int_{|t|<\delta}\bigl[f(x-t)-f(x)\bigr]\mathcal{F}_n(t)\,dt\biggr|
\le\frac\varepsilon\pi\int_{|t|<\delta}\mathcal{F}_n(t)\,dt
\le\frac\varepsilon\pi\int_{-\pi}^{\pi}\mathcal{F}_n(t)\,dt=\varepsilon.
\]

\textbf{第二部分} $\delta\le|t|\le\pi$。设 $M=\|f\|_\infty$，则 $|f(x-t)-f(x)|\le2M$。此时分母满足 $\sin^2\frac{t}{2}\ge\sin^2\frac\delta2$，故
\[
\mathcal{F}_n(t)=\frac{1}{2(n+1)}\frac{\sin^2\bigl(\frac{n+1}{2}t\bigr)}{\sin^2\frac{t}{2}}
\le\frac{1}{2(n+1)\sin^2\frac\delta2}.
\]
因此
\[
\frac1\pi\biggl|\int_{\delta\le|t|\le\pi}\bigl[f(x-t)-f(x)\bigr]\mathcal{F}_n(t)\,dt\biggr|
\le\frac{2M}{\pi}\int_{\delta\le|t|\le\pi}\mathcal{F}_n(t)\,dt
\le\frac{2M}{\pi}\cdot 2\pi\cdot\frac{1}{2(n+1)\sin^2\frac\delta2}
=\frac{2M}{(n+1)\sin^2\frac\delta2}.
\]

综合两部分，对任意 $x$ 一致地有
\[
|\sigma_n(f,x)-f(x)|
<\varepsilon+\frac{2M}{(n+1)\sin^2\frac\delta2}.
\]
取 $N$ 足够大，使得当 $n>N$ 时 $\frac{2M}{(n+1)\sin^2\frac\delta2}<\varepsilon$，则
\[
|\sigma_n(f,x)-f(x)|<2\varepsilon,
\]
故 $\sigma_n(f,x)$ 在 $[-\pi,\pi]$ 上一致收敛于 $f(x)$。由周期性，收敛在整个实轴上一致。
\end{proof}

以上证明充分体现了线性正算子与核的非负性在逼近论中的核心作用，与 Bernstein 多项式的证明如出一辙。


%\subsection{半无穷区间上的推广：Baskakov 与 Mirakyan 算子}

\subsubsection{Baskakov 线性正算子序列}

研究完闭区间 $[0,1]$ 上的 Weierstrass 第一逼近定理和 Bernstein 多项式后，一个自然的想法是：如果函数定义在半无穷区间 $[0, +\infty)$ 上，我们是否还能构造类似的函数列来一致逼近？这里单纯使用多项式是做不到的，需要引入一些更复杂的构造。

Baskakov 在 1957 年提出了如下线性正算子序列：
\[
\mathcal{B}_n(f;x) = \sum_{k=0}^{\infty} f\left( \frac{k}{n} \right) \binom{n+k-1}{k} x^k (1+x)^{-n-k},
\]
其中 $x \in [0, +\infty)$。这个构造的直观动机依然与概率论密切相关。在 Bernstein 多项式中，我们利用的是二项分布；而这里的核函数 $b_{n,k}(x) = \binom{n+k-1}{k} x^k (1+x)^{-n-k}$，本质上是取成功概率 $p = \frac{1}{1+x}$ 后的\textbf{负二项分布}。这意味着我们仍然在对函数 $f$ 进行某种“加权平均”，只不过求和是在无穷级数上进行的。

为了证明该算子一致收敛到 $f(x)$，我们需要建立类似于 Bernstein 多项式情形下的组合恒等式。首先，考虑广义二项式定理：当 $|z| < 1$ 时，
\[
(1-z)^{-n} = \sum_{k=0}^{\infty} \binom{n+k-1}{k} z^k.
\]
对这一展开式可作如下理解：
\[
(1-z)^{-n} = \left( \sum_{m=0}^{\infty} z^m \right)^n = \sum_{k=0}^{\infty} \binom{n+k-1}{k} z^k.
\]
基于此，我们通过对 $z$ 求导建立三个关键引理。

\begin{lemma}
\label{lem:baskakov1}
\[
\sum_{k=0}^{\infty} \binom{n+k-1}{k} x^k (1+x)^{-n-k} = 1.
\]
\end{lemma}
\begin{proof}
令 $z = \frac{x}{1+x}$（当 $x \geq 0$ 时 $0 \leq z < 1$），由广义二项式定理有 $\sum_{k=0}^{\infty} \binom{n+k-1}{k} z^k = (1-z)^{-n}$。两边同乘 $(1+x)^{-n}$ 并代回 $z$ 即得。该引理表明算子具有单位性。
\end{proof}

\begin{lemma}
\label{lem:baskakov2}
\[
\sum_{k=0}^{\infty} k \binom{n+k-1}{k} x^k (1+x)^{-n-k} = nx.
\]
\end{lemma}
\begin{proof}
对广义二项式定理两边关于 $z$ 求导，得 $n(1-z)^{-n-1} = \sum_{k=0}^{\infty} k \binom{n+k-1}{k} z^{k-1}$。两边同乘 $z(1+x)^{-n}$ 并代入 $z = \frac{x}{1+x}$ 即得。
\end{proof}

\begin{lemma}
\label{lem:baskakov3}
\[
\sum_{k=0}^{\infty} k^2 \binom{n+k-1}{k} x^k (1+x)^{-n-k} = n(n+1)x^2 + nx.
\]
\end{lemma}
\begin{proof}
对引理 \ref{lem:baskakov2} 对应的 $z$ 恒等式再次关于 $z$ 求导并同乘 $z$，或利用二阶导数性质即可导出。
\end{proof}

有了上述工具，我们可以仿照 Weierstrass 定理的证明思路，估计算子与 $f(x)$ 之间的偏差。

首先，利用引理 \ref{lem:baskakov1}--\ref{lem:baskakov3}，得到如下组合恒等式：
\begin{align*}
\sum_{k=0}^{\infty} (k-nx)^2 b_{n,k}(x)
&= \sum_{k=0}^{\infty} k^2 b_{n,k}(x) - 2nx \sum_{k=0}^{\infty} k b_{n,k}(x) + n^2x^2 \sum_{k=0}^{\infty} b_{n,k}(x) \\
&= [n(n+1)x^2 + nx] - 2nx(nx) + n^2x^2 \\
&= nx(1+x).
\end{align*}

下面考虑在任意给定的闭区间 $[0, a]$ 上的一致收敛性。设 $f(x)$ 在该区间上连续且有界（$|f(x)| \leq M$）。由一致连续性，对任给 $\epsilon > 0$，存在 $\delta > 0$，使得当 $|k/n - x| < \delta$ 时，有 $|f(k/n) - f(x)| < \epsilon$。

将求和按 $|k/n - x| < \delta$ 和 $|k/n - x| \geq \delta$ 分为两部分进行估计：
\begin{align*}
\left| \sum_{k=0}^{\infty} f\left(\frac{k}{n}\right) b_{n,k}(x) - f(x) \right|
&= \left| \sum_{k=0}^{\infty} [f(k/n) - f(x)] b_{n,k}(x) \right| \\
&\leq \sum_{|k/n - x| < \delta} \left|[f(k/n) - f(x)] b_{n,k}(x) \right| + \sum_{|k/n - x| \geq \delta} \left|[f(k/n) - f(x)] b_{n,k}(x) \right| \\
&< \epsilon \sum_{k=0}^{\infty} b_{n,k}(x) + \frac{2M}{n^2\delta^2} \sum_{k=0}^{\infty} (k-nx)^2 b_{n,k}(x) \\
&= \epsilon + \frac{2M \cdot nx(1+x)}{n^2\delta^2} = \epsilon + \frac{2Mx(1+x)}{n\delta^2}.
\end{align*}

在区间 $[0, a]$ 上，$x(1+x) \leq a(1+a)$ 有界。因此对于选定的 $\epsilon$ 和 $\delta$，只要 $n$ 足够大，使得 $n > \frac{2Ma(1+a)}{\epsilon\delta^2}$，偏差就能控制在 $2\epsilon$ 以内。这与 Bernstein 多项式的情形非常类似，本质上都是通过分段估计来控制误差。当然，也可以用 Korovkin 定理来处理这一问题。

综上所述，算子序列 $\mathcal{B}_n(f;x)$ 在 $[0, a]$ 上一致收敛到 $f(x)$。

\subsubsection{Mirakyan 算子序列}

设 $f$ 是 $\mathbb{R}$ 上的连续函数。定义 Mirakyan 算子序列
\[
M_n(f;x) = \sum_{k=0}^{\infty} f\left(\frac{k}{n}\right) \frac{n^k}{k!} x^k e^{-nx},
\]
下面证明 $M_n(f;x)$ 在任意闭区间 $[0, a]$ 上一致收敛于 $f$。

核函数 $\dfrac{(nx)^k}{k!}e^{-nx}$ 即为 Poisson 分布取参数 $\lambda = nx$ 的概率质量函数。Poisson 分布与二项分布关系密切：在 Weierstrass 第一逼近定理中，我们通过二项分布构造了 Bernstein 多项式 $B_n^f(x) = \sum_{k=0}^{n} f(\frac{k}{n}) \binom{n}{k} x^k (1-x)^{n-k}$，而这里相当于将核函数替换为 Poisson 分布。注意到 Poisson 极限公式
\[
\lim_{m \to \infty} \binom{m}{k} \left(\frac{nx}{m}\right)^k \left(1 - \frac{nx}{m}\right)^{m-k} = \frac{(nx)^k}{k!} e^{-nx},
\]
因此 Mirakyan 算子的收敛性也可以用来证明 Weierstrass 第一逼近定理。

由 $f$ 在 $[0,a]$ 上的一致连续性，$\forall \epsilon >0,\exists \delta>0$，使得当 $|x-y|<\delta$ 且 $x,y \in [0,a]$ 时，$|f(x)-f(y)|<\epsilon$。与 Weierstrass 定理的证明类似，这里同样有一个基本恒等式：
\[
\sum_{k=0}^{\infty} \frac{n^k}{k!} x^k e^{-nx} = 1.
\]
这由指数函数的 Taylor 展开可直接验证：
\[
\sum_{k=0}^{\infty} \frac{n^k}{k!} x^k e^{-nx}
= e^{-nx} \sum_{k=0}^{\infty} \frac{(nx)^k}{k!}
= e^{-nx} \cdot e^{nx} = 1.
\]

令 $|M_n(f;x)-f(x)| = \left|\sum_{k=0}^{\infty} \left(f\left(\frac{k}{n}\right) - f(x)\right)\frac{n^k}{k!} x^k e^{-nx}\right|$，则
\begin{align*}
|M_n(f;x)-f(x)|
&= \left| \left(\sum_{|k/n - x| < \delta} + \sum_{|k/n - x| \geq \delta}\right) \left(f\left(\frac{k}{n}\right) - f(x)\right) \frac{n^k}{k!} x^k e^{-nx} \right|.
\end{align*}

对于第一部分，利用一致连续性，
\begin{align*}
\left| \sum_{|k/n - x| < \delta} \left(f\left(\frac{k}{n}\right) - f(x)\right) \frac{n^k}{k!} x^k e^{-nx} \right|
&\leq \sum_{|k/n - x| < \delta} \left| f\left(\frac{k}{n}\right) - f(x) \right| \frac{n^k}{k!} x^k e^{-nx} \\
&\leq \sum_{k=0}^{\infty} \epsilon \frac{n^k}{k!} x^k e^{-nx}
= \epsilon \sum_{k=0}^{\infty} \frac{n^k}{k!} x^k e^{-nx} = \epsilon.
\end{align*}

对于第二部分，由于 $f$ 在 $[0,a]$ 上连续故有界，设 $\|f\| \leq M$，则 $|f(\frac{k}{n}) - f(x)| \leq 2M$。于是
\begin{align*}
\left| \sum_{|k/n - x| \geq \delta} \left(f\left(\frac{k}{n}\right) - f(x)\right) \frac{n^k}{k!} x^k e^{-nx} \right|
&\leq \sum_{|k/n - x| \geq \delta} 2M \frac{n^k}{k!} x^k e^{-nx}.
\end{align*}

由条件 $|\frac{k-nx}{n}| \geq \delta$ 可得 $(k-nx)^2 \geq (n\delta)^2$，因此
\begin{align*}
\sum_{|k/n - x| \geq \delta} 2M \frac{n^k}{k!} x^k e^{-nx}
&\leq \sum_{k=0}^{\infty} \frac{(k-nx)^2}{(n\delta)^2} 2M \frac{n^k}{k!} x^k e^{-nx} \\
&= \frac{2M}{n^2\delta^2} \sum_{k=0}^{\infty} (k-nx)^2 \frac{n^k}{k!} x^k e^{-nx}.
\end{align*}

为计算上述求和，利用指数函数幂级数的性质：
\[
\sum_{k=0}^{\infty} \frac{x^k}{k!} = e^x,
\]
\[
\sum_{k=0}^{\infty} k \cdot \frac{x^k}{k!}
= \sum_{k=1}^{\infty} \frac{x^k}{(k-1)!}
= \sum_{l=0}^{\infty} \frac{x^{l+1}}{l!}
= x e^{x},
\]
\[
\sum_{k=0}^{\infty} k^2 \frac{x^k}{k!}
= \sum_{k=1}^{\infty} k \cdot \frac{x^k}{(k-1)!}
= \sum_{k=1}^{\infty} (k-1) \cdot \frac{x^k}{(k-1)!} + \sum_{k=1}^{\infty} \frac{x^k}{(k-1)!}
= x^2 e^x + x e^x.
\]

代入可得
\begin{align*}
\frac{2M}{n^2\delta^2} \sum_{k=0}^{\infty} (k-nx)^2 \frac{n^k}{k!} x^k e^{-nx}
&= \frac{2M}{n^2\delta^2} \sum_{k=0}^{\infty} (k^2 - 2nxk + n^2x^2) \frac{n^k}{k!} x^k e^{-nx} \\
&= \frac{2M}{n^2\delta^2} \left( (nx)^2 e^{nx} + (nx) e^{nx} - 2nx(nx) e^{nx} + (nx)^2 e^{nx} \right) e^{-nx} \\
&= \frac{2Mx}{n\delta^2}.
\end{align*}

取 $N = \left[\frac{2Ma}{\epsilon \delta^2}\right]+1$，则当 $n > N$ 时，上式小于 $\epsilon$。综合两部分估计，得到
\[
|M_n(f;x)-f(x)| < 2\epsilon
\]
对任意 $x \in [0,a]$ 成立。因此 $M_n(f;x)$ 在 $[0,a]$ 上一致收敛于 $f(x)$。


%\input{chapter3_korovkin/theorem/vallee_poussin.tex}

\subsection{Korovkin 定理——线性正算子收敛的统一判据}

前面几节我们逐一给出了 Bernstein、Baskakov、Mirakyan 以及 V-P 算子的构造与收敛性证明。回顾这些证明，不难发现它们共享着几乎完全相同的论证结构：利用核函数的非负性与归一性，将逼近误差按远近拆分为两部分，近处用一致连续性控制，远处用函数的界与核函数的方差型求和控制，最终令 $n\to\infty$ 得到一致收敛。

这种高度雷同的证明模式暗示着背后存在一个统一的原理。Korovkin 在 20 世纪 50 年代发现了这一原理：对于线性正算子，其一致收敛性完全由三个“测试函数”的逼近行为所决定。在代数情形下，测试函数为 $1$, $x$, $x^2$；在三角情形下，测试函数为 $1$, $\cos x$, $\sin x$。本节将阐述并证明这两个 Korovkin 定理，并展示如何用它们简洁地统一处理前述所有算子。

\subsubsection{Korovkin 定理（代数多项式情形）}

\begin{theorem}[Korovkin 定理，代数多项式情形]
设线性正算子 $L_n: C[a,b]\to C[a,b]$ 满足
\[
L_n(1;x)=1+\alpha_n(x),\quad L_n(t;x)=x+\beta_n(x),\quad L_n(t^2;x)=x^2+\theta_n(x),
\]
其中 $\alpha_n,\beta_n,\theta_n$ 在 $[a,b]$ 上一致收敛于 $0$。则对任意 $f\in C[a,b]$，$L_n(f;x)$ 一致收敛于 $f(x)$。
\end{theorem}

\begin{proof}
由 $f$ 在 $[a,b]$ 上一致连续且 $\|f\|\le M$，$\forall \epsilon>0,\exists\delta>0$，当 $|x-y|<\delta$ 时 $|f(x)-f(y)|<\epsilon$。对任意 $x,t\in[a,b]$ 有
\[
-\frac{2M}{\delta^2}(t-x)^2 - \epsilon \le f(t)-f(x) \le \frac{2M}{\delta^2}(t-x)^2 + \epsilon.
\]
作用线性正算子 $L_n$（关于变量 $t$，$x$ 视为参数），利用线性性与非负性保持不等式方向，得
\[
-\frac{2M}{\delta^2}L_n((t-x)^2;x)-\epsilon L_n(1;x)
\le L_n(f;x)-f(x)L_n(1;x)
\le \frac{2M}{\delta^2}L_n((t-x)^2;x)+\epsilon L_n(1;x).
\]
计算
\begin{align*}
L_n((t-x)^2;x)
&= L_n(t^2;x) - 2xL_n(t;x) + x^2L_n(1;x)\\
&= (x^2+\theta_n(x)) - 2x(x+\beta_n(x)) + x^2(1+\alpha_n(x))\\
&= \theta_n(x) - 2x\beta_n(x) + x^2\alpha_n(x),
\end{align*}
它一致收敛于 $0$。又 $L_n(1;x)=1+\alpha_n(x)\to 1$。所以
\[
L_n(f;x) - f(x) = \big(L_n(f;x)-f(x)L_n(1;x)\big) + f(x)\alpha_n(x) \to 0
\]
一致成立。
\end{proof}

\subsubsection{Korovkin 定理（三角多项式情形）}

对于周期函数，有完全平行的结论。

\begin{theorem}[Korovkin 定理，三角多项式情形]
设线性正算子 $L_n: C_{2\pi}\to C_{2\pi}$ 满足
\[
L_n(1;x)=1+\alpha_n(x),\quad L_n(\cos t;x)=\cos x+\beta_n(x),\quad L_n(\sin t;x)=\sin x+\theta_n(x),
\]
且 $\alpha_n,\beta_n,\theta_n$ 一致收敛于 $0$，则对任意 $f\in C_{2\pi}$，$L_n(f;x)$ 一致收敛于 $f(x)$。
\end{theorem}

\begin{proof}
证明思路与代数情形类似。由 $f$ 在 $\R$ 上一致连续且以 $2\pi$ 为周期，$\forall\epsilon>0,\exists\delta>0$，当 $|x-y|<\delta$ 时 $|f(x)-f(y)|<\epsilon$。对任意 $x,t\in\R$，可用三角不等式得到
\[
-\frac{2M}{\sin^2(\delta/2)}\sin^2\frac{t-x}{2} - \epsilon \le f(t)-f(x) \le \frac{2M}{\sin^2(\delta/2)}\sin^2\frac{t-x}{2} + \epsilon.
\]
利用恒等式
\[
\sin^2\frac{t-x}{2} = \frac{1}{2}\big(1 - \cos t\cos x - \sin t\sin x\big),
\]
对 $L_n$ 作用后，$L_n(\sin^2\frac{t-x}{2};x)$ 可展开为 $L_n(1;x)$, $L_n(\cos t;x)$, $L_n(\sin t;x)$ 的线性组合。由定理条件，这三项分别一致收敛于 $1$, $\cos x$, $\sin x$，因此 $L_n(\sin^2\frac{t-x}{2};x)$ 一致收敛于 $0$。其余步骤与代数情形完全相同。
\end{proof}

\subsubsection{统一验证：回到经典算子}

Korovkin 定理的美妙之处在于，它将原本需要逐一进行的分段估计，统一简化为对三个测试函数的验证。下面以 Bernstein 多项式和 Fejér 算子为例，展示这种统一性在代数情形与三角情形中的平行应用。

\textbf{Bernstein 多项式}：对 $f\in C[0,1]$，$B_n(f;x)=\sum_{k=0}^{n} f(\frac{k}{n})\binom{n}{k}x^{k}(1-x)^{n-k}$ 是线性正算子。直接计算三个测试函数：
\[
B_n(1;x)=1,\qquad B_n(t;x)=x,\qquad B_n(t^{2};x)=x^{2}+\frac{x(1-x)}{n}.
\]
因此 $\alpha_n(x)=0$，$\beta_n(x)=0$，$\theta_n(x)=\frac{x(1-x)}{n}\rightrightarrows 0$。由 Korovkin 定理即得 $B_n(f;x)$ 一致收敛于 $f(x)$。

\textbf{Fejér 算子}：对周期函数 $f\in C_{2\pi}$，Fejér 算子
\[
\sigma_n(f,x)=\frac{1}{\pi}\int_{-\pi}^{\pi}f(x-t)\mathcal{F}_n(t)\,dt,\qquad
\mathcal{F}_n(t)=\frac{1}{2(n+1)}\frac{\sin^{2}\bigl(\frac{n+1}{2}t\bigr)}{\sin^{2}\frac{t}{2}}
\]
是线性正算子（积分核非负）。利用三角形式的 Korovkin 定理，只需验证测试函数 $1,\cos t,\sin t$。

由归一化条件即得 $\sigma_n(1,x)=1$。对于 $\cos t$，利用 $\cos(x-t)=\cos x\cos t+\sin x\sin t$ 以及 $\mathcal{F}_n$ 为偶函数，有
\[
\sigma_n(\cos t,x)=\frac{1}{\pi}\int_{-\pi}^{\pi}\cos(x-t)\mathcal{F}_n(t)\,dt
=\cos x\cdot\frac{1}{\pi}\int_{-\pi}^{\pi}\cos t\,\mathcal{F}_n(t)\,dt.
\]
由 Dirichlet 核的性质，$\frac{1}{\pi}\int_{-\pi}^{\pi}\cos t\,D_k(t)\,dt=1$（$k\ge 1$）而 $k=0$ 时为 $0$，因此
\[
\frac{1}{\pi}\int_{-\pi}^{\pi}\cos t\,\mathcal{F}_n(t)\,dt
=\frac{1}{n+1}\sum_{k=0}^{n}\frac{1}{\pi}\int_{-\pi}^{\pi}\cos t\,D_k(t)\,dt
=\frac{n}{n+1}=1-\frac{1}{n+1}.
\]
于是 $\sigma_n(\cos t,x)=\frac{n}{n+1}\cos x=\cos x-\frac{1}{n+1}\cos x$，余项一致趋于 $0$。同理，由奇偶性可得
\[
\sigma_n(\sin t,x)=\sin x\cdot\frac{1}{\pi}\int_{-\pi}^{\pi}\cos t\,\mathcal{F}_n(t)\,dt
=\frac{n}{n+1}\sin x=\sin x-\frac{1}{n+1}\sin x,
\]
余项亦一致趋于 $0$。由三角形式 Korovkin 定理立即得到 $\sigma_n(f,x)\rightrightarrows f(x)$，这样 Weierstrass 第二逼近定理作为 Korovkin 定理的直接推论而成立，与 Bernstein 多项式的情形形成完美对照。

\subsection{小结}

本章从具体的构造性算子出发，给出了 Weierstrass 两个逼近定理的完整证明，随后引入 Korovkin 定理，揭示了这些算子收敛性的共同本质：只需验证三个测试函数的矩条件，即可保证对一切连续函数的一致收敛。Bernstein 多项式的二项分布核、Baskakov 算子的负二项分布核、Mirakyan 算子的 Poisson 分布核、Vallée-Poussin 算子的幂次余弦核，虽然形式各异，但在 Korovkin 的框架下，它们的收敛性验证都归结为完全相同的简单计算。

然而，Korovkin 定理只回答了“是否收敛”的问题，并未涉及逼近的速率。要刻画“收敛有多快”，需要将线性算子的误差与最佳逼近误差联系起来，这正是下一章 Lebesgue 不等式的任务。